% Jiri Kotek xkotek09
% VUT FIT 2023

\documentclass [11pt, a4paper, twocolumn] {article} 

\usepackage[utf8]{inputenc}
\usepackage{times}
\usepackage[czech]{babel}
\usepackage[T1]{fontenc}
\usepackage[left=1.4cm, text={18.2cm, 25.2cm}, top=2.3cm]{geometry}
\usepackage{hyperref}
\usepackage{hyphenat}
\usepackage{amsthm}
\usepackage{amsmath}
\usepackage{icomma}
\usepackage{ dsfont }
\usepackage{ stmaryrd }
\usepackage{enumitem}
\newtheorem{name}{Printed output}
\newtheorem{Definice}{Definice}
\newtheorem{name1}{Printed output1}
\newtheorem{Veta}{Veta}

\setlist[itemize]{itemsep=0.6ex, topsep=0.6ex}

\title{Typografie a publikování  \,--\, 2. projekt}
\author{Jiří Kotek \\ \href{mailto:xkotek09@stud.fit.vutbr.cz}{xkotek09@stud.fit.vutbr.cz}}
\date{}


\begin{document}
    
    \begin{titlepage}
        \begin{center}
        \Huge
        \textsc{Vysoké učení technické v Brně\\
        \huge
        {Fakulta informačních technologií}\\}
        \vspace{\stretch{0.382}}
        \LARGE
        Typografie a publikování\,--\,2. projekt \\
        Sazba dokumentů a matematických výrazů
        \vspace{\stretch{0.618}}
        \end{center} 
        {\Large 2023 \hfill
        Jiří Kotek(xkotek09)}
    \end{titlepage}

    \section*{Úvod}
V této úloze si vyzkoušíme sazbu titulní strany, matematických vzorců, prostředí a dalších textových struktur obvyklých pro technicky zaměřené texty \,--\, například Definice 1 nebo rovnice (\ref{3}) na straně \pageref{3}. 
Pro vytvoření těchto odkazů používáme kombinace příkazů \verb|\label|, \verb|\ref|, \verb|\eqref| a \verb|\pageref|. 
Před odkazy patří nezlomitelná mezera. Pro zvýrazňování textu jsou zde několikrát použity příkazy \verb|\verb| a \verb|\emph|.

Na titulní straně je použito prostředí \verb|titlepage| a sázení nadpisu podle optického středu s využitím přesného zlatého řezu. 
Tento postup byl probírán na přednášce. 
Dále jsou na titulní straně použity čtyři různé velikosti písma a~mezi dvojicemi řádků textu je použito odřádkování se zadanou relativní velikostí 0${,}$5 em a 0${,}$4 em\footnote [1]{Nezapomeňte použít správný typ mezery mezi číslem a jednotkou}.

\section{Matematický text}
V této sekci se podíváme na sázení matematických symbolů a výrazů v plynulém textu pomocí prostředí \verb|math|.
Definice a věty sázíme pomocí příkazu \verb|\newtheorem| s~využitím balíku \verb|amsthm|. Někdy je vhodné použít konstrukci~\verb|${}$|~nebo~\verb|\mbox{}|,~která~říká,~že~\mbox{(matematický)} text nemá být zalomen. 

\begin{Definice}
    ${}$ \textnormal{Zásobníkový automat} (ZA) je definován jako sedmice tvaru $A = (Q, \Sigma, \Gamma, \delta, q_{0}, Z_{0}, F) $, kde:

    \begin{itemize}
        \item $Q$  je konečná množina \textnormal{vnitřních (řídících) stavů,}
        \item $\Sigma $ je konečná \textnormal{vstupní abeceda,}
        \item $\Gamma $ je konečná \textnormal{zásobníková abeceda,}
        \item $\delta$   je \textnormal{přechodová funkce} $Q \times (\Sigma \cup \{\epsilon\}) \times \Gamma \rightarrow 2^{Q \times \Gamma^*} $,
        \item $q_{0} \in Q $  je \textnormal{počáteční stav, $Z_{0} \in \Gamma $ }je\textnormal{ startovací symbol zásobníku a $F \subseteq Q $ } je množina\textnormal{ koncových stavů.}
    \end{itemize}
\end{Definice}
 
Nechť $P = (Q, \Sigma, \Gamma, \delta, q_{0}, Z_{0}, F) $ je ZA. \textit{Konfigurací} nazveme trojici $(q, w, \alpha ) \in Q \times \Sigma^* \times \Gamma^* $, kde $q$ je aktuální stav vnitřního řízení, $w$ je dosud nezpracovaná část vstupního řetězce a $\alpha = Z_{i_{1}} Z_{i_{2}} \dots Z_{i_{k}}$ je obsah zásobníku.

\subsection{Podsekce obsahující definici a větu}

\begin{Definice}
    \textnormal{Řetězec} $ w $ \textnormal{nad abecedou} $ \Sigma $ \textnormal{je přijat ZA}~$ A $~jest\-liže \mbox{$(q_{0}, w, Z_{0}) \overset{*}{\underset{A}{\vdash}} (q_{F}, \epsilon, \gamma) $} pro nějaké \mbox{$\gamma \in \Gamma^{*}  $} a \mbox{$q_{F} \in F $}.
    Množina $L(A) = \{w \;|\; w$ je přijat ZA A$\} \subseteq \Sigma^{*} $ je \textnormal{jazyk příjímaný} ZA A.
\end{Definice}

\begin{Veta}
    Třída jazyků, které jsou příjímány ZA, odpovídá \textnormal{bezkontextovým jazykům.}
\end{Veta}

\section{Rovnice}
Složitější matematické formulace sázíme mimo plynulý text pomocí prostředí \verb|displaymath|. 
Lze umístit i několik výrazů na jeden řádek, ale pak je třeba tyto vhodně oddělit, například příkazem \verb|\quad|.
\begin{displaymath}
    1^{2^{3}} \neq \Delta_{\Delta_{\Delta^{3}}^{2}}^{1}
    \quad
    y_{22}^{11}-\sqrt[9]{x + \sqrt[7]{y}}
    \quad
    x > y_{1} \leq y^{2}
\end{displaymath}
V rovnici (\ref{2}) jsou využity tři typy závorek s různou \textit{ex\-plicitně} definovanou velikostí. 
Také nepřehlédněte, že nasledující tři rovnice mají zarovnaná rovnítka, a použijte k~tomuto účelu vhodné prostředí. 
\begin{alignat}{2}
    -\cos^{2} \beta &= \frac{\frac{\frac{1}{x}+\frac{1}{3} }{y}+1000 }{ \overset{8}{\underset{j=2}{\prod}} q_{j} }\\
    \bigg( \bigl\{ b \star \big[ 3 \div 4 \big] \circ a \bigr\}^{\frac{2}{3}}    \bigg) &= \log_{10} x \label{2}\\
    \int_a^b  f(x)\, \mathrm{d}x &= \int_c^d  f(y)\, \mathrm{d}y \label{3}
\end{alignat}
V této větě vidíme, jak vypadá implicitní vysázení limity $\lim_{m \rightarrow \infty} f(m)  $ v normálním odstavci textu. 
Podobně je to i s dalšími symboly jako $\bigcup_{N \in \mathcal{M}} N $ či  $ \sum_{i=1}^{m} x_{i}^{2} $. S vynucením méně úsporné sazby příkazem \verb|\limits| budou vzorce vysázeny v podobě $ \lim\limits_{m \rightarrow \infty} f(m) $ a $ \sum\limits_{i=1}^{m} x_{i}^{4} $.

\section{Matice}
Pro sázení matic se velmi často používá prostředí \verb|array| a závorky(\verb|\left |, \verb|\right |) .
    \[
        \mathbf{B}=
    \left|
    \begin{array}{cccc} 
        b_{11} & b_{12} & \dots & b_{1n}\\
        b_{21} & b_{22} & \dots & b_{2n}\\
        \vdots & \vdots & \ddots & \vdots \\
        b_{m1} & b_{m2} & \dots & b_{mn}\\
    \end{array}
    \right|
    =
    \left|
    \begin{array}{cc} 
        t & u\\
        v & w\\
    \end{array}
    \right|
    =
    tw - uv
    \]

    \[
        \mathds{X} =  \mathbf{Y} \Longleftrightarrow 
    \begin{bmatrix}
         & \Omega  +\Delta & {\psi}  \\
        \overset{\shortrightarrow}{\pi} & \omega & \\
    
    \end{bmatrix}
    \neq 42 \]
    

Prostředí\,\verb|array|\,lze úspěšně využít i jinde, například~na pravé straně následující rovnice. 
Kombinační číslo na~levé straně vysázejte pomocí příkazu \verb|\binom|.

$$
\binom{n}{k} =
\left\{
    \begin{array}{ccc}
        0 & \textnormal{pro } k < 0\\
        \frac{n!}{k! (n - k)!} & \textnormal{pro } 0 \leq k \leq n \\
        0 & \textnormal{pro } k > 0 \\
    \end{array}
    \right.
$$

\end{document}